\subsection*{IV.17.15　勘误段（续）}
\begingroup
\def\OO{\mathcal O}\def\m{\mathfrak m}\def\mm{\mathfrak m}\def\Sym{\operatorname{Sym}}
反过来，IV.23.9.2 按原样是不对的。关于例外范围，法文稿写的是：$k$的特征不等于$2$，且$X$的各不可约分支维数均为偶数。\footnote{法文第113页清楚印作加下划线的“et”（且），而英译第74页作“or”（或）。此处标出法稿的写法，不把它当成已证的最宽范围；后面的独立校核仅在所查到的 Hessian 判据下确认“或”。这两点不等于断定原稿所指的整条 IV.23.9.2 在何种情形成立。}
在一般情形，应删去“特别地，$f_n$满足8.8中的等价条件”。命题的其余部分看来是正确的，应该可以用相当类似的方法来证明：对于给定的$x\in X(k)$和指定的次数$n\ge2$，证明存在一个次数为$n$的超曲面$H$，它在$x$处与$X$相切，使$x$成为$X\cap H$的一个初等二次点，且交截的维数少一……。

不妨把问题说透：若$k$的特征为$2$，而$X$非空、连通且为奇数维，则8.8的条件不成立，也就是说，$L/K$必定不可分。（我想若$X\subset\mathbb P^r$，次数恰为$2$，不过没有把握……。）\footnote{法文第114页将“非空”手写补入。英译所留法语惯用语直译会失去原意，此处译成直接指出问题；“没有把握”是作者对次数的保留，未将其改作定理。}
\paragraph{译校说明（非原文）。}
\paragraph{草稿中的交叉引用。}
“$f_n$满足8.8中的等价条件”与本英译预备笔记第36页命题9.2（改变射影嵌入）高度吻合：在$X$光滑、$f$非分歧的设置下，选取整数$n\ge2$，研究第$n$次 Veronese 复合$f_n=u_n\circ f$。第28页在有关泛点处 étale、有限可分扩张的论述里写作定理8.7(c)，而第30--31页又使用8.8。这里仅给出极可能的同稿线索，保留手稿编号，不认定它是排印版《EGA IV》中有同编号的命题。

\paragraph{与 Katz 的 Hessian 判据对照。}
参见 Katz，\href{https://grothendiecksga.com/read/sga7/fr/17-3.html}{《SGA 7 II》，第XVII讲，§§3.3.15、3.5.0、3.7.1}。在该讲采用的代数闭域、光滑不可约射影情形，相切关联概形投影在指定点非分歧的判据是 Hessian 可逆。特征$2$时 Hessian 为交替矩阵；若$\dim X$为奇数，则不可能可逆。对于普通二次奇点，在$\operatorname{char}k\ne2$\emph{或}$\dim X$为偶数时，光滑二次切锥的射影化确实保证极化非退化。这是上述析取在特定判据中的依据；它并不单独证明整条 EGA 草稿命题，更不补足跨编号的条件。

\paragraph{给定点的局部实现。}
在上述光滑射影设置，取$x=[1:0:\cdots:0]$并使$X_i/X_0$（$1\le i\le \dim X$）在$x$处给出参数。若$Q$是这些参数中的二次型，则齐次多项式$X_0^{n-2}Q(X_1,\ldots,X_{\dim X})$在$x$处常数、线性项均为零，二次初始项为$Q$。选择其射影零概形光滑的$Q$即可给出草稿所说的局部初等二次点。特征$2$且维数为奇数时，取若干双曲对加未配对变量的平方，仍可得到光滑射影零概形，但极化矩阵不可逆。这个计算不主张所选$H$在别处也光滑。

\paragraph{可检验的二次不可分实例。}
特征$2$下对$\mathbb P^1$取二次 Veronese 嵌入，切超平面对应的二元二次型在判别集$b=0$上可写成$aX^2+cY^2$。在$a\ne0$、坐标$t=X/Y$处，泛型切点满足$t^2=c/a$。参数$c/a$在判别集的函数域中不是平方（在零点处的赋值为$1$），因而该例的扩张确为纯不可分二次。它不把草稿关于一般$X$的猜测升级为结论，也不为任意情形的$L/K$提供一般次数证明。
\endgroup
